% ===========================================================================
%  user/02-overview/00-chapter.tex
% ===========================================================================
\chapter{Overview}\label{cap:overview}

It is the page that opens first. It answers a single question: how is
everything going, right now.

\figuraSchermata{1.0}{figure/schermate/overview.png}
  {The home screen. Four panels: the ensemble chart, the metrics, the list of
   active streams, and the shortcut to Studio.}
  {overview}

\section{The chart at the top left}\label{sez:overview-chart}

It shows the performance of the chosen \voce{Project}: the combination of
several strategies, net of estimated costs, on the data the strategy had
\emph{not} seen while it was being built.

Until a project has been saved, an invitation to create one takes its place.
It is not an error: there is simply nothing to draw yet.

\section{Metrics OOS}\label{sez:overview-metrics}

Six numbers, and they are meant to be read together. \voce{OOS} stands for
\emph{out of sample}: computed on data the strategy did not know.

\begin{center}
\begin{tabular}{@{}L{32mm}L{106mm}@{}}
\toprule
\textbf{Metric} & \textbf{What it tells you} \\
\midrule
\rigaTenue
\voce{Sharpe ratio} & How much it earns relative to how much it swings. It is
  the number most used to compare two strategies, and the easiest to make
  look good by looking at the wrong data. \\
\voce{CAGR net} & The compound annual return, already net of estimated
  costs. \\
\rigaTenue
\voce{Max drawdown} & The largest loss from a previous peak. It is the
  question \enquote{how much would I have had to endure before recovering}. \\
\voce{Ann. volatility} & How much it swings in a year. \\
\rigaTenue
\voce{Alpha vs SPY ann.} & How much it earns \emph{beyond} the market, once
  the part explained by the market itself is removed. A strategy with a good
  return and an alpha close to zero is buying the market under other
  names. \\
\voce{Win rate} & The percentage of days in profit. On its own it says
  little: you can win nine times out of ten and still lose. \\
\bottomrule
\end{tabular}
\end{center}

\begin{perche}{Which numbers to look at first}
\voce{Max drawdown} and \voce{Alpha vs SPY}. The first because it is the
figure that decides whether a strategy can be lived with: an excellent Sharpe
with a \perc{45} drawdown is a strategy nobody holds to the end. The
second because it tells a strategy that does something apart from one that
follows the market with longer strides.
\end{perche}

\begin{attenzione}{A backtest is not a forecast}
It says how a strategy \emph{would have} behaved on past data, with
estimated costs. It does not say how it will behave. \vediCap{streams}
explains which signs tell a strategy that worked from one that got lucky.
\end{attenzione}

\section{Active streams}\label{sez:overview-list}

The list of all available results, sorted. The columns:

\begin{description}[leftmargin = 1.4em, style = nextline, itemsep = 4pt]
  \item[\voce{Sharpe}] As above.
  \item[\voce{OOS Return}] The out-of-sample return. Green above zero, red
    below.
  \item[\voce{Corr}] The correlation with the other streams. It is the column
    that matters most when combining: two highly correlated strategies, put
    together, diversify nothing --- they make the same bet twice.
  \item[\voce{Watch}] The star adds the stream to your favourites.
\end{description}

\noindent
The three menus above the list choose the period (\voce{OOS} or all of it),
the sort order and the reference project.

\section{Tutorials}\label{sez:overview-tutorials}

At the bottom right, four steps that tick themselves off as you do them:
build the first stream, validate it, combine it, put it into production. It
is the path this guide follows in the next chapters.
